\scanpage{109}
% archive foliation: 108
\textcircled{a}\quad $(P,\sigma,c)$

$P$ fibré en droites projectives sur $X$

$c$ connexion sur $P$ rel$/S$

$\sigma$ section partout non stationnaire

\textcircled{$a'$}\quad triple $(P,\sigma,\varphi)$, $(P,\sigma)$ comme dessus,
$\varphi$ iso \deleted{$\unsure{\mathcal P}$} $\underline O_{\widehat{(P,\sigma)}}\simeq P^\infty_{X/S}$ d'Alg. augm. telle que la connexion transportée sur $\underline O_{\widehat{(P,\sigma)}}$ \deleted{\unsure{par}} se globalise à $P$

\textcircled{b}\quad Algèbre de Lie de type $\mathfrak{sl}(\unsure{2})$ quasi-scindée
\[
\mathcal L_u:=\Omega\subset\mathcal L\subset\underbrace{\mathcal H}_{\mathrm{Sym}^2(\mathcal L)\otimes\Omega^{-1}}
\qquad(\text{NB }\mathcal L/\Omega\simeq\underline O_X)
\]
munie de connexion $c\deleted{\unsure{D}}:\mathcal H\longrightarrow\mathcal H\otimes\underline\omega$
(\unsure{néc.t} $c(\Omega)\subset\mathcal L(1)$) telle que la composée
\[
\Omega\xrightarrow{c_\Omega}\mathcal L\otimes\omega\longrightarrow\omega
\]
soit un iso (\uline{NB} Il suffit de se donner $c$ sur $\mathcal L$ \deleted{\unsure{ce qui}} car
\[
\mathcal L\xrightarrow{c_{\unsure{\mathcal L|2}}}\mathrm{Sym}^2\mathcal L\otimes\Omega^{-1}\otimes\omega\simeq\mathrm{Sym}^2\mathcal L
\]
--- ``compatible au crochet'' i.e. ....)

\textcircled{$b'$}\quad Sur l'Algèbre de Lie \inserted{\uline{quasi-scindée} $\mathcal H_X$}
\[
\mathrm{Sym}^2(P^1_{X/S}(\underline\omega)\otimes\underline\omega^{-1})\otimes\underline\omega^{-1}
\simeq\mathrm{Sym}^2(P^1_{X/S}(\underline\omega))\otimes\underline\omega^{-3},
\]
\deleted{une connexion}
(avec $\mathcal L\unsure{\subset|=}P^1_{X/S}(\underline\omega)\otimes\underline\omega^{-1}\unsure{\}\,.}\supset\underline\omega^2\otimes\underline\omega^{-1}\simeq\underline\omega$)
une connexion $c$ d'Algèbre de Lie
\deleted{\unsure{qui sur $\omega$}} $c:\mathcal H_{\unsure{1|X}}\longrightarrow\mathcal H_{\unsure{1|X}}\otimes\underline\omega$
qui sur $\underline\omega$ soit $d^1_{X/S,\underline\omega}:\underline\omega\longrightarrow P^1_{X/S}(\underline\omega)$

\textcircled{c}\quad Donné $\mathbb P_X=\mathbb P(P^1_{X/S}(\underline\omega))$, munie de $\sigma_X$, et du \deleted{iso} can.
$\underline O_{\mathbb P_X}/\mathcal J^3\overset{\varphi_2}{\simeq}P^2_{X/S}$, un iso
$\underline O_{\mathbb P_X}/\mathcal J^4\overset{\varphi_3}{\simeq}P^3_{X/S}$ qui le prolonge

\textcircled{d}\quad Une équa. diff. \inserted{(du 3\textsuperscript{e} ordre rel. $X$)} sur $\mathbb P^1_S\times_S X\simeq\mathbb P^1_X$ \deleted{rel. à $X$}, (à ne pas confondre avec $\mathbb P_X$ !), invariante par $\unsure{\mathrm{GP}(1)_X|\mathrm{GPM}_X}$ opérant sur $\mathbb P^1_X$, --- équa. \unsure{par.} aux jets d'ordre 3 de \unsure{sections} i.e. morphismes $X\longrightarrow\mathbb P^1_S$ qui sont des iso locaux à l'ordre 3
